% ch7_e 第7章 §7.4尾 (书页 327–334 / p342–p349)
%JOIN%
% ——— 块边界说明（起点）———
% 书页 326（p341）末句为 "We find that in the new coordinates the components of
%     the gauge field"，跨页延续到书页 327（p342）顶部 "are given by" + 式 (7.4.34)。
% 上半句译文（"我们发现，在新坐标下规范场的各分量"）由 ch7_d 收尾；
% 本块正文以 %JOIN% 起头，从"由下式给出"+式 (7.4.34) 续起，与 ch7_d 末句缝合。
% ——— 块边界说明（终点）———
% 本块为章末块：书页 334（p349）止于问题 7.4.9，其后空白，第 7 章自然收束
% （书页 335 即第 8 章章首），故本块不设 %--D-TAIL-- 区。
% ——— 覆盖范围 ———
% §7.4.5 后部（自式 (7.4.34) 起，含式 (7.4.35)–(7.4.37)、问题 7.4.7、脚注 49）；
% §7.4.6 全节（式 (7.4.38)–(7.4.48)、问题 7.4.8、脚注 50）；
% §7.4.7 全节（图 7.18、表 7.3、式 (7.4.49)、问题 7.4.9）。
% ——————————————————————————————
由下式给出：
\begin{equation*}
\tilde{a}_{\tilde{t}}=a_{t}+va_{x},\qquad \tilde{a}_{\tilde{x}}=a_{x},\qquad \tilde{a}_{\tilde{y}}=a_{y}. \tag{7.4.34}
\end{equation*}
在新坐标下，规范固定条件就化为前面讨论过的那个条件。容易看出，在式 (7.4.33) 与式 (7.4.34) 给出的变换下，Chern--Simons 作用量的形式保持不变：
\begin{equation*}
S=-\frac{m}{4\pi}\int\mathrm{d}^{3}x\,a_{\mu}\partial_{\nu}a_{\lambda}\varepsilon^{\mu\nu\lambda}
=-\frac{m}{4\pi}\int\mathrm{d}^{3}x\,\tilde{a}_{\tilde{\mu}}\partial_{\tilde{\nu}}\tilde{a}_{\tilde{\lambda}}\varepsilon^{\tilde{\mu}\tilde{\nu}\tilde{\lambda}}
\end{equation*}
重复前面的推导，我们发现边缘作用量为
\begin{equation*}
S=-\frac{m}{4\pi}\int\mathrm{d}\tilde{t}\,\mathrm{d}\tilde{x}\,\partial_{\tilde{t}}\phi\,\partial_{\tilde{x}}\phi
\end{equation*}
用原来的物理坐标表示，上述作用量具有形式
\begin{equation*}
S=-\frac{m}{4\pi}\int\mathrm{d}t\,\mathrm{d}x\,(\partial_{t}+v\partial_{x})\phi\,\partial_{x}\phi \tag{7.4.35}
\end{equation*}
这是一个手征玻色子理论（chiral boson theory）(Gross et al., 1985; Floreanini and Jackiw, 1988)。由运动方程 $\partial_{t}\phi+v\partial_{x}\phi=0$ 可以看出，式 (7.4.35) 所描述的边缘激发具有非零的速度 $v$，并且只朝一个方向运动。

为了得到式 (7.4.35) 的量子理论，我们需要把手征玻色子理论量子化。量子化可以在动量空间中进行：
\begin{equation*}
S=\frac{m}{2\pi}\int\mathrm{d}t\sum_{k>0}\left(\mathrm{i}k\dot{\phi}_{k}\phi_{-k}-vk^{2}\phi_{k}\phi_{-k}\right).
\end{equation*}
可以证明，如果我们把 $k>0$ 的 $\phi_{k}$ 视为``坐标''，那么 $\phi_{-k}$ 将正比于与之对应的动量 $\pi_{k}=\partial S/\partial\dot{\phi}_{k}$。辨认出``坐标''与``动量''之后，我们便可以得到哈密顿量以及对易关系 $[\phi_{k},\pi_{k^{\prime}}]=\mathrm{i}\delta_{kk^{\prime}}$，它们完全定义了这个量子系统。如果引入 $\rho=\frac{1}{2\pi}\partial_{x}\phi$，我们发现该量子系统由下式描述：
\begin{equation*}
[\rho_{k},\rho_{k^{\prime}}]=\frac{k\delta_{k+k^{\prime}}}{2\pi m},\qquad
H=2\pi mv\sum_{k>0}\rho_{k}^{\dagger}\rho_{k} \tag{7.4.36}
\end{equation*}

边缘激发的速度 $v$ 通过规范固定条件进入我们的理论。注意，在受限规范变换下，具有不同 $v$ 的规范固定条件 (7.4.32) 彼此不能互相变换，它们在物理上并不等价。这与我们的假设一致：规范固定条件中的 $v$ 是一个物理参量，它实际上决定了边缘激发的速度。

哈密顿量 (7.4.36) 只有当 $vm<0$ 时才有下界。理论的自洽性要求 $v$ 与 $m$ 符号相反。因此，边缘激发速度的符号（即手性（chirality））由 Chern--Simons 项前面系数的符号决定。

由于矩阵 $K$ 可以对角化，上述结果可以容易地推广到式 (7.3.18) 所描述的一般 FQH 态。所得到的有效边缘理论具有形式
\begin{equation*}
S_{\mathrm{edge}}=\frac{1}{4\pi}\int\mathrm{d}t\,\mathrm{d}x\,\bigl[K_{IJ}\partial_{t}\phi_{I}\partial_{x}\phi_{J}-V_{IJ}\partial_{x}\phi_{I}\partial_{x}\phi_{J}\bigr] \tag{7.4.37}
\end{equation*}
哈密顿量为
\begin{equation*}
H_{\mathrm{edge}}=\frac{1}{4\pi}\int\mathrm{d}t\,\mathrm{d}x\,V_{IJ}\partial_{x}\phi_{I}\partial_{x}\phi_{J}
\end{equation*}
因此，$V$ 必须是正定矩阵。利用这个结果，可以证明 $K$ 的正本征值对应于一个左移分支，而负本征值对应于一个右移分支。

$\nu=2/5$ FQH 态的有效理论由 $K=\begin{pmatrix}3 & 2\\ 2 & 3\end{pmatrix}$ 给出。由于 $K$ 具有两个正本征值，$\nu=2/5$ FQH 态的边缘激发有两个朝同一方向运动的分支。$\nu=1-\frac{1}{n}$ FQH 态则由 $K=\begin{pmatrix}1 & 0\\ 0 & -n\end{pmatrix}$ 的有效理论描述。此时 $K$ 的两个本征值符号相反，因此边缘激发的两个分支朝相反的方向运动。

\subsection*{问题 7.4.7}
证明：量子化之后，系统 (7.4.35) 由式 (7.4.36) 描述。

\subsection{带电激发与电子传播子}

\begin{keypoints}
\item 一般 FQH 边缘上的一般电子/准粒子算符。
\end{keypoints}

我们已经研究了几种简单 FQH 液体中边缘激发的动力学，发现低能边缘激发由一个自由声子理论描述。本节我们将集中讨论一般的电荷激发。特别地，我们将对最一般的（阿贝尔）FQH 态计算电子与准粒子的传播子。关键点同样是把电子或准粒子算符用声子算符 $\rho_{I}$ 写出来。一旦做到这一点，传播子便可以容易地算出，因为在低能与长波极限下声子是自由的。

我们知道，对于式 (7.3.18) 所描述的 FQH 态，其边缘态由作用量 (7.4.37) 描述。边缘激发的希尔伯特空间构成 K--M 代数的一个表示
\begin{align*}
[\rho_{Ik},\rho_{Jk^{\prime}}]&=(K^{-1})_{IJ}\frac{1}{2\pi}k\,\delta_{k+k^{\prime}}\\
k,k^{\prime}&=\text{整数}\times\frac{2\pi}{L} \tag{7.4.38}
\end{align*}
其中 $\rho_{I}=\frac{1}{2\pi}\partial_{x}\phi_{I}$ 是 FQH 态中第 $I$ 个凝聚体的边缘密度，$I,J=1,\ldots,\kappa$，而 $\kappa$ 是 $K$ 的维数。边缘上的电子密度为
\begin{equation*}
\rho_{e}=-eq_{I}\rho_{I}
\end{equation*}
边缘激发的动力学由哈密顿量
\begin{equation*}
H=2\pi\sum_{IJ}V_{IJ}\rho_{I,-k}\rho_{J,k} \tag{7.4.39}
\end{equation*}
描述，其中 $V_{IJ}$ 是正定矩阵。

我们先尝试写出边缘上的准粒子算符 $\Psi_{\pmb{l}}$，它产生一个由 $l_{I}$ 标记的准粒子。我们知道，把准粒子插到边缘上会使第 $I$ 个凝聚体的边缘密度发生一个改变 $\delta\rho_{I}$（见式 (7.3.22)）。$\delta\rho_{I}$ 满足
\begin{equation*}
\int\mathrm{d}x\,\delta\rho_{I}=(\pmb{l}^{\top}K^{-1})_{I}.
\end{equation*}
由于 $\Psi_{\pmb{l}}$ 是一个只引起密度局域改变的局域算符，我们有
\begin{equation*}
[\rho_{I}(x),\Psi_{\pmb{l}}(x^{\prime})]=l_{J}(K^{-1})_{JI}\delta(x-x^{\prime})\Psi_{\pmb{l}}(x^{\prime}) \tag{7.4.40}
\end{equation*}
利用 Kac--Moody 代数 (7.4.38)，可以证明满足式 (7.4.40) 的准粒子算符为
\begin{equation*}
\Psi_{\pmb{l}}\propto\mathrm{e}^{\mathrm{i}\phi_{I}l_{I}} \tag{7.4.41}
\end{equation*}
准粒子 $\Psi_{\pmb{l}}$ 的电荷由对易子 $[\rho_{e},\Psi_{\pmb{l}}]$ 确定，结果为
\begin{equation*}
Q_{\pmb{l}}=-e\,\pmb{q}^{\top}K^{-1}\pmb{l} \tag{7.4.42}
\end{equation*}
由式 (7.3.24) 与式 (7.4.41)，我们看到电子算符可以写成
\begin{equation*}
\Psi_{e,L}\propto\mathrm{e}^{\mathrm{i}\sum_{I}l_{I}\phi_{I}},\qquad
l_{I}=K_{I,J}L_{J},\qquad q_{I}L_{I}=1 \tag{7.4.43}
\end{equation*}
由式 (7.4.42) 可见，上述算符携带单位电荷。$\Psi_{e,L}$ 的对易可如下计算：
\begin{align*}
\Psi_{e,L}(x)\Psi_{e,L}(x^{\prime})&=(-)^{\lambda}\Psi_{e,L}(x^{\prime})\Psi_{e,L}(x)\\
\lambda&=l_{I}(K^{-1})_{IJ}l_{J}=L_{I}K_{IJ}L_{J} \tag{7.4.44}
\end{align*}
在分级基（hierarchical basis）中，$K_{II}|_{I=1}$ 为奇数，$K_{II}|_{I>1}$ 为偶数，$\pmb{q}^{\top}=(1,0,\ldots,0)$，且 $L_{1}=1$。利用这些条件，我们可以证明 $(-)^{\lambda}=-1$。式 (7.4.43) 定义的电子算符的确是费米子算符。

由于不同 $L_{I}$ 选择下的所有算符 $\Psi_{e,L}$ 都携带单位电荷并且是费米型的，每一个 $\Psi_{e,L}$ 都可能是电子算符的候选。一般地，真正的电子算符是诸 $\Psi_{e,L}$ 的下述叠加：
\begin{equation*}
\Psi_{e}=\sum_{L}C_{L}\Psi_{e,L}
\end{equation*}
这里，当我们说边缘上存在许多不同的电子算符时，我们的真正意思是：真实的物理电子算符是这些算符的一个叠加。

利用 K--M 代数 (7.4.38) 与哈密顿量 (7.4.39)，我们可以计算一个一般准粒子算符
\begin{equation*}
\Psi_{\pmb{l}}\propto\mathrm{e}^{\mathrm{i}l_{I}\phi_{I}}
\end{equation*}
（对于 $\pmb{l}$ 的合适选择，其中包括电子算符）的传播子。首先，我们注意到变换
\begin{equation*}
\rho_{I}\rightarrow U_{IJ}\tilde{\rho}_{J}
\end{equation*}
把 $(V,K)$ 变为 $(\tilde{V},\tilde{K})=(U^{\top}VU,\,U^{\top}KU)$。选取合适的 $U$，$K$ 与 $V$ 可以被同时对角化；\footnote{由于 $V$ 是正定的对称矩阵，我们可以找到一个 $U_{1}$，它把 $V$ 变为 $V_{1}=U_{1}^{\top}VU_{1}=1$，把 $K$ 变为 $K_{1}=U_{1}^{\top}KU_{1}$。注意 $K_{1}$ 是一个对称矩阵，其本征值与 $K$ 的本征值同号（尽管它们的绝对值可以不同）。然后我们做一个正交变换把 $K_{1}$ 对角化，即 $(K_{1})_{IJ}\rightarrow(K_{2})_{IJ}=\sigma_{I}|v_{I}|^{-1}\delta_{IJ}$。在这个正交变换下，$V_{1}\rightarrow V_{2}=1$ 不变且保持对角。最后，我们用一个对角矩阵把 $K_{2}$ 缩放为 $(K_{3})_{IJ}=\sigma_{I}\delta_{IJ}$。这把 $V_{2}$ 变为 $V_{3}=|v_{I}|\delta_{IJ}$。新的 $(V_{3},K_{3})$ 便导致式 (7.4.45)。}用 $\tilde{\rho}_{I}$ 表示，式 (7.4.38) 与式 (7.4.39) 变为
\begin{align*}
[\tilde{\rho}_{Ik},\tilde{\rho}_{Jk^{\prime}}]&=\sigma_{I}\delta_{IJ}\frac{1}{2\pi}k\,\delta_{k+k^{\prime}}\\
H&=2\pi\sum_{I,k>0}|v_{I}|\,\tilde{\rho}_{I,-k}\tilde{\rho}_{I,k} \tag{7.4.45}
\end{align*}
其中 $\sigma_{I}=\pm1$ 是 $K$ 本征值的符号。由 $\tilde{\rho}_{I}$ 产生的边缘激发的速度为 $v_{I}=\sigma_{I}|v_{I}|$。

用 $\tilde{\rho}_{I}$ 表示，算符 $\Psi_{\pmb{l}}$ 具有形式
\begin{equation*}
\Psi_{\pmb{l}}\propto\mathrm{e}^{\mathrm{i}\sum_{I}\tilde{l}_{I}\tilde{\phi}_{I}},\qquad
\tilde{l}_{J}=l_{I}U_{IJ} \tag{7.4.46}
\end{equation*}
由式 (7.4.45) 与式 (7.4.46)，我们看到 $\Psi_{\pmb{l}}$ 的传播子具有以下一般形式：
\begin{equation*}
\langle\Psi_{\pmb{l}}^{\dagger}(x,t)\Psi_{\pmb{l}}(0)\rangle\propto
\mathrm{e}^{\mathrm{i}l_{I}k_{I}x}\prod_{I}(x-v_{I}t+\mathrm{i}\sigma_{I}\delta)^{-\gamma_{I}},\qquad
\gamma_{I}=\tilde{l}_{I}^{2} \tag{7.4.47}
\end{equation*}
其中 $v_{I}=\sigma_{I}|v_{I}|$ 是边缘激发的速度。右移成分的总指数减去左移成分的总指数，满足求和规则\footnote{我们用到了
$\lambda_{\pmb{l}}=\sum_{I,J,I^{\prime}}l_{I^{\prime}}U_{I^{\prime}I}\sigma_{I}(U^{\top})_{I,J}l_{J}$，以及
$(K^{-1})_{IJ}=U_{II^{\prime}}\sigma_{I^{\prime}}\delta_{I^{\prime},J^{\prime}}(U^{\top})_{J^{\prime},J}$。}
\begin{equation*}
\sum_{I}\sigma_{I}\gamma_{I}\equiv\lambda_{\pmb{l}}=\pmb{l}^{\top}K^{-1}\pmb{l} \tag{7.4.48}
\end{equation*}
我们看到，差值 $\lambda_{\pmb{l}}$ 不依赖边缘相互作用 $V$，是一个拓扑量子数。由式 (7.4.44) 与式 (7.4.48)，我们看到这个差值与 $\Psi_{\pmb{l}}$ 的统计性质直接相关。如果 $\Psi_{\pmb{l}}$ 代表电子算符（或其他费米型算符），那么 $\lambda_{\pmb{l}}$ 将是一个奇整数。从式 (7.4.47) 我们还看到，算符 $\Psi_{\pmb{l}}$ 产生的激发带有接近 $\sum_{I}l_{I}k_{I}$ 的动量。

\subsection*{问题 7.4.8}
证明式 (7.4.41) 中的 $\Psi_{\pmb{l}}$ 满足式 (7.4.40)。

\subsection{手征 Luttinger 液体的唯象后果}

\begin{keypoints}
\item 三类边缘态。
\item 长程相互作用与杂质散射的效应。
\end{keypoints}

在前四节中，我们已经证明 FQH 液体边缘上的电子构成手征 Luttinger 液体。作为手征 Luttinger 液体的特征性质之一，电子与准粒子的传播子获得如下的反常指数（anomalous exponents）：
\begin{equation*}
\langle\Psi_{e}^{\dagger}(t,x=0)\Psi_{e}(0)\rangle\sim t^{-g_{e}},\qquad
\langle\Psi_{q}^{\dagger}(t,x=0)\Psi_{q}(0)\rangle\sim t^{-g_{q}}
\end{equation*}
这些反常指数可以通过 FQH 边缘之间的隧穿实验直接测量。

%FIG{7.18}: {(a) 同一 FQH 态的两条边缘之间的准粒子隧穿。(b) 不同 FQH 态的边缘之间的电子隧穿。}

以下两种情形需要分开考虑：(i) 两条 FQH 液体边缘被一个绝缘体隔开；(ii) 两条边缘被一个 FQH 液体隔开。情形 (i) 中，只有电子能够在两条边缘之间隧穿（见图 7.18(b)）；情形 (ii) 中，隔开两条边缘的 FQH 液体所支持的准粒子能够隧穿（见图 7.18(a)）。情形 (i) 的隧穿算符为 $A\propto\Psi_{e1}\Psi_{e2}^{\dagger}$，其中 $\Psi_{e1}$ 与 $\Psi_{e2}$ 是两条边缘上的电子算符；情形 (ii) 的隧穿算符具有形式 $A\propto\Psi_{q1}\Psi_{q2}^{\dagger}$，其中 $\Psi_{q1}$ 与 $\Psi_{q2}$ 是两条边缘上的准粒子算符。隧穿的物理性质可以由隧穿算符的关联函数算出，而后者又可以表示为两条边缘上电子或准粒子传播子的乘积。

令 $g=g_{e}$ 对应情形 (i)，$g=g_{q}$ 对应情形 (ii)。那么在零温下，电子与准粒子算符的反常指数导致 $\langle A(t)A(0)\rangle\propto|t|^{-2g}$，从而给出一条非线性的隧穿 $I$--$V$ 曲线（见 3.7.4 节）
\begin{equation*}
I\propto V^{2g-1}
\end{equation*}
隧穿电流的噪声谱在频率 $f=QV/h$ 处还包含一个奇异性，即
\begin{equation*}
S(f)\sim\bigl|f-\frac{QV}{h}\bigr|^{2g-1} \tag{7.4.49}
\end{equation*}
其中 $V$ 是两条边缘之间的电压差，$Q$ 是隧穿电子或隧穿准粒子的电荷。在有限温度 $T$ 下，零偏压电导也具有如下幂律依赖关系：
\begin{equation*}
\sigma\propto T^{2g-2}
\end{equation*}
我们看到，反常指数可以很容易地通过隧穿实验测量。噪声谱还进一步揭示了隧穿粒子的电荷。

指数 $g_{e}$ 与 $g_{q}$ 已在 7.4.6 节中算出。为了以简单的方式总结这些结果，最好把 FQH 边缘分成以下两类：(A) 所有边缘激发朝同一方向运动；(B) 边缘激发朝两个方向传播。

对于 A 类边缘，指数 $g_{e}$ 与 $g_{q}$ 与电子及准粒子的统计性质直接相关。在这种情形下，$g_{e}$ 与 $g_{q}$ 是普适的，不依赖电子相互作用和边缘势的细节。表 7.3 列出了一些支持 A 类边缘态的 FQH 态，以及相应的指数 $g_{e,q}$ 的数值和相应粒子的电荷（注意表中只列出了 $g_{e}$ 与 $g_{q}$ 的最小值）。

\begin{center}
\small 表 7.3\quad 具有 A 类边缘态的某些 FQH 态的 $(g_{e},g_{q},Q_{q})$\\[6pt]
\begin{tabular}{cccccc}
\toprule
FQH 态 & $\nu=1/m$ & $\nu=2/5$ & $\nu=p/(pq+1)$ & $(331)_{\nu=1/2}$ & $(332)_{\nu=2/5}$ \\
\midrule
$g_{e}$ & $m$ & $3$ & $q+1$ & $3$ & $3$ \\
$g_{q}$ & $1/m$ & $2/5$ & $p/(pq+1)$ & $3/8$ & $2/5$ \\
电荷 $Q_{q}/e$ & $1/m$ & $2/5$ & $p/(pq+1)$ & $1/4$ & $2/5$ \\
\bottomrule
\end{tabular}
\end{center}

让我们讨论一下，对于填充因子为 $\nu=\frac{p}{pq+1}$（$q=$ 偶数）的分级态，上述结果是如何得到的。这些态包括 $\nu=2/5,3/7,2/9,\ldots$ 等 FQH 态。填充因子为 $\nu=\frac{p}{pq+1}$ 的分级态，在 $\pmb{q}^{\top}=(1,1,\ldots,1)$ 的基下由 $p\times p$ 矩阵 $K=1+qC$ 描述。这里 $C$ 是伪单位矩阵（pseudo-identity matrix），即 $C_{IJ}=1$，$I,J=1,\ldots,p$，而 $K^{-1}=1-\frac{q}{pq+1}C$。由于所有边缘激发都朝同一方向运动，我们有
\begin{equation*}
\langle\Psi_{\pmb{l}}^{\dagger}(x=0,t)\Psi_{\pmb{l}}(0)\rangle\propto t^{-\lambda_{\pmb{l}}}
\end{equation*}
其中 $\lambda_{\pmb{l}}$ 由式 (7.4.48) 给出。基本准粒子由 $\pmb{l}^{\top}=(1,0,\ldots,0)$ 给出，携带电荷 $\frac{1}{pq+1}$，其传播子中的指数为 $\lambda_{\pmb{l}}=1-\frac{q}{pq+1}$。指数最小的准粒子由 $\pmb{l}^{\top}=(1,\ldots,1)$ 给出，携带电荷 $\frac{p}{pq+1}$，指数为 $\lambda_{\pmb{l}}=\frac{p}{pq+1}$，它小于 $1-\frac{q}{pq+1}$（注意我们有 $q\geqslant2$ 与 $p\geqslant1$）。在低能下，正是这样的准粒子（携带电荷 $\frac{p}{pq+1}$）主导\emph{同一} FQH 流体两条边缘之间的隧穿（见 7.4.7 节）。

电子算符由 $\Psi_{e,L}=\Psi_{\pmb{l}}$（译注：原文排印为 $\psi_{\pmb{l}}$，应为 $\Psi_{\pmb{l}}$ 之误）给出，其中 $\pmb{l}$ 满足 $\sum_{I}l_{I}=pq+1$。传播子中的指数为 $\lambda_{\pmb{l}}=\sum l_{I}^{2}-q(pq+1)$。传播子中指数最小的电子算符由 $\pmb{l}^{\top}=(q,\ldots,q,q+1)$ 给出，最小指数的值为 $\lambda_{\pmb{l}}=q+1$。在低能下，正是这样的电子算符主导两个\emph{不同} FQH 流体边缘之间的隧穿。

对于 B 类边缘，$g_{e}$ 与 $g_{q}$ 不是普适的。它们的数值依赖于电子相互作用和边缘势。

\subsection*{问题 7.4.9}
利用有限电压 $V$ 下隧穿算符的关联 $\langle A(t)A(0)\rangle$，推导式 (7.4.49)。
